\section{问题提出的背景}

\subsection{背景介绍}

多模态大模型的发展经历了从视觉-语言对齐到统一多模态推理，再到视频时空建模的逐步演进。CLIP 通过大规模图文对比学习验证了视觉特征与语言语义统一建模的有效性，为后续多模态模型提供了表示基础\cite{clip2021}；Flamingo、BLIP-2 和 LLaVA 等工作进一步通过跨模态注意力、桥接模块与视觉指令微调等机制，形成了“视觉编码器-投影器-大语言模型”的主流架构范式\cite{flamingo2022,blip22023,llava2023}。在此基础上，研究开始从静态图像扩展到动态视频，Video-LLaMA、Video-LLaVA 与 VideoLLaMA2 相继引入音频-视频-文本联合建模、统一视觉表示以及 STC 时空连接器，使视频大模型具备更强的长时序理解与跨模态推理能力\cite{videollama2023,videollava2023,videollama22024}。

然而，能力提升也同步放大了推理代价。对于视频输入，帧数增加会直接带来视觉 token 数量膨胀；跨模态融合又会放大注意力计算与显存访存压力；在长视频场景下，时序特征缓存和上下文建模进一步成为实时部署的主要障碍\cite{llamavid2024,longvlm2024}。近年的 PruneVid 等研究表明，视频大模型的视觉 token 中存在显著冗余，若能在保持关键信息的前提下进行有效压缩，能够明显改善推理效率\cite{prunevid2024}。因此，本课题聚焦“多模态视频大模型高效推理”这一问题，选择在 VideoLLaMA2-7B-16F 的视觉编码器与 STC connector 之间引入 TokenPacker 压缩模块，探索在尽量不改动主干结构的前提下，以逐帧空间压缩缓解后续时空融合与语言解码阶段的计算负担\cite{tokenpacker2024,videollama22024}。

本课题题目为《多模态视频大模型高效推理关键技术研究》，核心研究目标为采用时空联合压缩技术，有效缓解视频大模型推理效率瓶颈，明确研究方向为将 TokenPacker 中的 token 压缩技术迁移到 VideoLLaMA2-7B-16F 模型，具体插入位置为“视觉编码器之后、STC connector 之前”。

\subsection{本研究的意义和目的}

本课题题目为《多模态视频大模型高效推理关键技术研究》，核心研究目标为采用时空联合压缩技术，有效缓解视频大模型推理效率瓶颈，研究方向为将 TokenPacker 中的 token 压缩技术迁移到 VideoLLaMA2-7B-16F 模型，具体插入位置为视觉编码器之后、STC connector 之前。

本研究的理论意义在于验证图像多模态中的细节感知型 token 压缩方法能否迁移到视频大模型场景。现有视频大模型的效率优化多聚焦于关键帧选择、时空下采样或 token 剪枝，但对于如何在压缩空间 token 的同时保持局部细节表达，仍缺少兼具明确机制与可复用实现的方案。TokenPacker 提出的“低分辨率粗表示+高分辨率多层细节注入”机制，为视频场景中的逐帧空间压缩提供了较强的方法启发。

本研究的实践意义在于，为 VideoLLaMA2-7B-16F 提供一种轻量、模块化的高效推理改造方式。在保持视觉编码器、STC connector 和 LLM 主体框架基本不变的条件下，通过在 vision tower 与 STC 之间插入 TokenPacker，实现对每帧 patch token 的先压缩再融合，有望减少后续时空聚合和语言解码阶段的计算负担，从而降低推理时延和显存开销。

本研究的具体目的包括：第一，完成 TokenPacker 在 VideoLLaMA2-7B-16F 中的结构接入与接口适配；第二，构建适用于视频场景的“逐帧空间压缩+STC 时空融合”推理流程；第三，通过对比实验分析不同压缩率下的视觉 token 数量、推理耗时、显存占用及视频理解性能变化；第四，总结 TokenPacker 迁移到视频大模型的有效条件、局限性及进一步优化方向。
